% TOP_PROBLEM: {"id": 3313, "title": "Every 4-connected toroidal graph has a Hamilton cycle", "category_id": 3, "category": {"id": 3, "name": "graph_theory", "display_name": "Graph Theory", "description": "Problems involving graphs, networks, and their properties.", "slug": "graph-theory", "order_index": 3, "created_at": "2026-07-31T15:26:25.671Z"}, "status": "open", "problem_number": "OPG-46943", "set_id": 12}
% FIRST_POSTED: 2026-10-02
% ATTEMPT_STATUS: unresolved
% MODEL: GPT 6 Astra Ultra
% UnsolvedMath ID 3313; source catalogue code OPG-46943
% !TeX program = lualatex
% Research attempt dated 2026-10-02; canonical statement preserved.
\documentclass[11pt,a4paper]{article}
\usepackage[margin=25mm]{geometry}
\usepackage{fontspec}
\setmainfont{lmroman10-regular.otf}[BoldFont=lmroman10-bold.otf,
 ItalicFont=lmroman10-italic.otf,BoldItalicFont=lmroman10-bolditalic.otf]
\usepackage{amsmath,amssymb,mathtools}
\usepackage{unicode-math}
\setmathfont{latinmodern-math.otf}
\AtBeginDocument{\let\setminus\smallsetminus}
\usepackage{xurl,hyperref}
\hypersetup{colorlinks=true,urlcolor=blue}
\setcounter{secnumdepth}{0}
\setlength{\parindent}{0pt}
\setlength{\parskip}{5pt}
\setlength{\emergencystretch}{3em}
\begin{document}
\section*{Every 4-connected toroidal graph has a Hamilton cycle}\label{3313}
{\small UnsolvedMath ID 3313 · OPG-46943 · Graph Theory · OpenGarden}
\subsection{Definitions and mathematical statement}
A finite simple undirected graph is toroidal if it can be embedded
without crossings in the torus $\mathbb T^2=S^1\times S^1$. In an
embedding, vertices are distinct points and edges are arcs meeting
only at common endpoints. Toroidal means admitting such an embedding;
it includes planar graphs and does not require the torus to be the
surface of minimum possible genus.

The graph $G$ is 4-connected if it has at least five vertices and,
for every $S\subseteq V(G)$ with $|S|\le3$, the graph $G-S$ remains
connected. This is vertex connectivity, not edge connectivity. A
Hamilton cycle is a simple cycle containing every vertex exactly
once, apart from repeating its starting vertex at the end.

The conjecture asks whether every finite simple 4-connected toroidal
graph contains a Hamilton cycle. Equivalently, if $V(G)$ has $n$
vertices, there should be an ordering $v_1,\ldots,v_n$ of all of them
such that
\[
                  v_iv_{i+1}\in E(G)\ (1\le i<n),
                  \qquad v_nv_1\in E(G).
\]
The cycle is not required to pass through a prescribed edge or to
represent a particular homotopy class on the torus. Its existence is
an abstract graph property, independent of which toroidal embedding
is chosen. The target is a closed cycle, so a spanning path with no
closing edge does not suffice. No triangulation assumption or
minimum-degree condition beyond what 4-connectivity implies is imposed.
\subsection{Short English statement}
Does every 4-connected graph that can be embedded on a torus have a cycle visiting every vertex?
\subsection{Research attempt}
\paragraph{Scope and method.}
This research attempt by GPT-6 Astra Ultra is dated 2 October 2026.
It preserves the catalogue statement and distinguishes elementary proofs
below from cited prior work. No novelty claim or formal proof-assistant
verification is made. The final paragraph states the precise remaining scope.
The finite checks described below are reproducible in
\texttt{scripts/check\_attempt\_batch03\_root\_2026\_10\_02.py}.
\paragraph{Target and literature boundary.}
The desired object is a spanning cycle. A spanning path does not suffice:
its two endpoints need not be adjacent. Thomas, Yu and Zang [R1] prove
the Hamilton-path theorem for 4-connected toroidal graphs. Ellingham
and Marshall [R2] discuss the stronger demand that every specified edge
belong to a Hamilton cycle and counterexamples to that strengthening.
Neither replacement target can silently be substituted for this one.
Here we construct spanning cycles in a concrete infinite family and prove
that the family meets the connectivity hypothesis.

\paragraph{A toroidal family and its connectivity.}
Let $G=C_m\mathbin\square C_n$, where $m,n\ge3$, with vertices
$(i,j)$ modulo $(m,n)$. Horizontal and vertical consecutive pairs are
edges. The rectangular drawing with opposite sides identified embeds it
on the torus. Each vertex has four distinct neighbours, even when a
factor has length three.
Delete a set $S$ of at most three vertices. Suppose first that $m\ge4$.
There is an intact row. Every column containing at most one deletion
is connected after deletion, and hence connects all its surviving
vertices to that row. At most one column can contain two or more
deletions. A surviving vertex in that column has two horizontal
neighbours in distinct other columns. At most one vertex outside the
exceptional column was deleted, so one of these neighbours survives.
It belongs to a good column. Thus every remaining vertex reaches the
intact row, proving connectivity. If $m=3<n$, interchange the factors.
For $m=n=3$ the same argument applies when a row or column is intact.
Otherwise the three deletions occupy distinct rows and columns. Since
each factor is $K_3$, relabel them as the diagonal. The six remaining
vertices form the cycle
$(0,1),(0,2),(1,2),(1,0),(2,0),(2,1),(0,1)$.
This proves 4-connectivity in all cases; deleting the four neighbours
of any vertex also shows that the connectivity is exactly four.

\paragraph{Constructive cycle splicing.}
Initially take the $m$ disjoint horizontal row cycles. For
$i=0,\ldots,m-2$, choose the horizontal edge on columns $0,1$ if
$i$ is even, and on columns $1,2$ if $i$ is odd. Delete that edge
from row $i$ and its corresponding edge from row $i+1$. Join the
corresponding endpoints by the two vertical edges between those rows.
At this stage the first deleted edge belongs to the cycle already
spanning rows $0,\ldots,i$, while the second belongs to a separate
row cycle. Cutting one edge from each and joining the resulting paths
crosswise produces one cycle on their combined vertex sets.
The edge required in row $i$ is still present: the preceding splice,
if any, used the other horizontal edge. They are distinct even for
$n=3$, although they may share a vertex. The induction therefore
continues until a Hamilton cycle on all $mn$ vertices remains.
No parity condition on either factor is necessary.

\paragraph{Why the construction does not prove the conjecture.}
A general toroidal embedding need not contain disjoint spanning row
cycles or the paired vertical edges used in a splice. Four-connectivity
prevents a small vertex separator; it does not supply this product
structure. Likewise, cutting an arbitrary essential curve on the
torus changes the boundary conditions of a prospective cycle. Rejoining
the cut does not automatically preserve a single spanning component.
A general induction would need an extension lemma controlling both
connectivity and how a cycle meets that boundary. No such lemma is
proved here, and the prescribed-edge strengthening cannot be used
as an unproved induction hypothesis.

\paragraph{Verification and final status.}
The accompanying root check script verifies the splice construction for
$3\le m,n\le12$ and checks every deletion of at most three vertices
for $3\le m,n\le6$. The parameter proofs above cover all sizes;
these finite checks only test their implementation. The restricted
product-family theorem is established, while the general toroidal
Hamilton-cycle question remains unresolved by this attempt.

\subsection{Sources}
\begin{itemize}
\item[S1] ULAM.AI and the UnsolvedMath Contributors, supplied problems.json, record 3313 (OPG-46943), OpenGarden collection. Catalogue identity and source transcription; CC BY 4.0.
\url{https://huggingface.co/datasets/ulamai/UnsolvedMath}.
\item[S2] Open Problem Garden, Every 4-connected toroidal graph has a Hamilton cycle. Original problem and discussion; the mathematical formulation below is expanded from the supplied source record.
\url{http://www.openproblemgarden.org/op/every_4_connected_toroidal_graph_has_a_hamilton_cycle}.

\item[R1] R. Thomas, X. Yu and W. Zang, Hamilton paths in toroidal graphs,
J. Combin. Theory B 94 (2005), 214--236. Author-institution record:
\url{https://hub.hku.hk/handle/10722/156133}.
\item[R2] M. N. Ellingham and E. A. Marshall, Criticality of counterexamples
to toroidal edge-hamiltonicity, arXiv:1312.1379.
\url{https://arxiv.org/abs/1312.1379}.

\end{itemize}
\end{document}
